\documentclass[11pt,a4paper]{article}
\usepackage[margin=1in]{geometry}
\usepackage{hyperref}
\usepackage{enumitem}

\title{Strategic Management \\ Week 10 \\ Slide-by-Slide Notes (ELI5) \\ Unit 1}
\author{(Generated from \texttt{sm-week-10-slides-2.pdf})}
\date{}

\begin{document}
\maketitle

\section*{How to use this (ELI5)}
\begin{itemize}
  \item For each slide, write the ``Exam write'' line in your notebook.
  \item Then copy the ELI5 explanation as your simple understanding.
  \item If you do not understand one slide, re-read that slide section only.
\end{itemize}

% ---------------------------
\section*{Slide 1 of 17 (Contents)}
\textbf{Key idea (from slide):} Unit 1 Strategic Management (course map / overview).\\
\textbf{ELI5:} This unit is about what strategy really means.\\
\textbf{Exam write:} ``Unit 1 explains what strategy is, the levels of strategy, and how strategic management works.''

% ---------------------------
\section*{Slide 2 of 17 (Learning objectives)}
\textbf{Key idea (from slide):}
\begin{itemize}
  \item Understand what strategic decisions mean.
  \item Identify critical elements of a successful strategy.
  \item Understand interrelations among strategy levels and business functions.
  \item Understand rational and emergent processes of strategic management.
\end{itemize}
\textbf{ELI5:} By the end, you should know the main parts of strategy and how they connect.\\
\textbf{Exam write:} ``My strategy answer must connect meaning, elements, levels, and process (rational + emergent).'' 

% ---------------------------
\section*{Slide 3 of 17 (Origins of strategy)}
\textbf{Key idea (from slide):} Slide introduces the origin idea of the word/idea of strategy.\\
\textbf{ELI5:} Strategy comes from war thinking first.\\
\textbf{Exam write:} ``Strategy has a military origin (plan + leadership), then became a business concept.''

% ---------------------------
\section*{Slide 4 of 17 (Epistemology of strategy)}
\textbf{Key idea (from slide):}
\begin{itemize}
  \item Military origins: ``stratós'' (army) + ``ēgós'' (leadership).
  \item Business administration use developed in academic context (Von Neumann and Morgenstern, 1944).
  \item Spread in the 1960s via Harvard Business School and BCG.
  \item Publications: Porter (1980, 1985).
\end{itemize}
\textbf{ELI5:} People used war ideas (leadership of an army) and later used them for business competition.\\
\textbf{Exam write:} ``Strategy = managing competition like a general manages a campaign under uncertainty.''

% ---------------------------
\section*{Slide 5 of 17 (What is strategy?)}
\textbf{Key idea (from slide):}
\begin{itemize}
  \item Success in a sustained way (no luck).
  \item Strategic decisions imply: uncertainty, complex decisions, expensive mistakes (even survival).
  \item Chandler definition: long-run goals/objectives + courses of action + resource allocation.
\end{itemize}
\textbf{ELI5:} Strategy is long-term planning for a future you cannot fully predict, and you cannot afford many mistakes.\\
\textbf{Exam write:} ``Strategy is long-run direction + actions + resources, chosen under uncertainty with costly mistakes.''

% ---------------------------
\section*{Slide 6 of 17 (What’s strategy? (Firm + Environment)) }
\textbf{Key idea (from slide):}
\begin{itemize}
  \item Strategy links THE FIRM (goals/values, resources/capabilities, structure, people) with THE ENVIRONMENT (competitors, customers, suppliers, national institutions).
\end{itemize}
\textbf{ELI5:} Strategy must fit two things: what your company is like, and what the market is like.\\
\textbf{Exam write:} ``In an answer, show fit between firm (resources/people/structure) and environment (competitors/customers/suppliers/institutions).''

% ---------------------------
\section*{Slide 7 of 17 (Definition of strategy: key elements)}
\textbf{Key idea (from slide):} Successful strategy needs:
\begin{itemize}
  \item Simple, consistent, long-term goals.
  \item Understanding competitive environment.
  \item Appraisal of internal resources.
  \item Effective implementation.
\end{itemize}
\textbf{ELI5:} Strategy is: decide long-term goal, understand the market, check your internal power, then execute well.\\
\textbf{Exam write:} ``A complete strategy has goals + external understanding + internal appraisal + implementation.''

% ---------------------------
\section*{Slide 8 of 17 (The Art of War)}
\textbf{Key idea (from slide):}
\begin{itemize}
  \item Sun Tzu treatise (360 BC).
  \item Goal: win at war.
  \item Apply Art of War ideas to four elements of successful strategy (using Grant, 2015).
\end{itemize}
\textbf{ELI5:} Use war logic to remember strategy elements: goal, market, internal strength, and execution.\\
\textbf{Exam write:} ``Use the Art of War idea to structure: long-run goals, environment reading, internal appraisal, and implementation.''

% ---------------------------
\section*{Slide 9 of 17 (Levels of strategy)}
\textbf{Key idea (from slide):} Three major levels:
\begin{itemize}
  \item Corporate strategy: set scope (activities/businesses), vision as unit of competition, firm-environment fit, value-added opportunities (internationalization, M\&A, alliances, diversification).
  \item Functional strategy: resource utilization by functional areas (ops, HR, marketing, financial, R\&D…).
  \item Competitive strategy: how to compete in a business; sustained competitive advantage; cost leadership / product differentiation.
  \item Multi-business firms use Strategic Business Units (SBU).
\end{itemize}
\textbf{ELI5:} Corporate = where you play, competitive = how you win in each game, functional = how each department supports winning.\\
\textbf{Exam write:} ``Always separate corporate vs competitive vs functional logic in answers.''

% ---------------------------
\section*{Slide 10 of 17 (Rational / deliberated process)}
\textbf{Key idea (from slide):} Deliberated-rational process:
\begin{itemize}
  \item Strategy formulation: external/internal analysis.
  \item Strategy implementation: organizational support, design, strategic management control and review.
  \item Internal analysis methods include SWOT; adequacy/reliability/acceptability checks.
\item External/internal analysis feeds corporate \& competitive strategies.
  \item Missions, vision, objectives, values appear inside the process.
\end{itemize}
\textbf{ELI5:} In a rational process, you diagnose, then choose, then implement, then check if it worked.\\
\textbf{Exam write:} ``State the process chain: analysis -> formulation -> implementation -> control/review.''

% ---------------------------
\section*{Slide 11 of 17 (Emergent strategy process + why rational is not enough)}
\textbf{Key idea (from slide):}
\begin{itemize}
  \item Advantages: managerial explicit process; clear and quantitative objectives/evaluation criteria.
\item But real decisions depend on limited rationality \& personal features of key actors (directors, CEO, TMT).
  \item Also depend on features of strategy: uncertainty, complexity, paradoxes, and learning process.
  \item Changing internal organization and environment.
  \item End includes political process and even luck-intuition.
\end{itemize}
\textbf{ELI5:} Real strategy is not a perfect math plan. People, learning, politics, and surprises change decisions.\\
\textbf{Exam write:} ``Explain emergent strategy: limited rationality + actors + learning + changing context + politics/luck effects.''

% ---------------------------
\section*{Slide 12 of 17 (Integrative process conclusion)}
\textbf{Key idea (from slide):}
\begin{itemize}
  \item Strategic management is intended + emergent.
  \item Subject to personal cognitive restrictions, desires/preferences.
  \item Subject to environmental dynamism (change).
\end{itemize}
\textbf{ELI5:} Strategy is both planned and adjusted.\\
\textbf{Exam write:} ``In answers: use integrative logic = planned intentions + emergent adjustment due to people and changing context.''

% ---------------------------
\section*{Slide 13 of 17 (Responsibilities: directors, CEO, TMT, staff)}
\textbf{Key idea (from slide):}
\begin{itemize}
  \item Board of Directors, CEO-president, heads of operations/marketing/funding, and heads by country.
  \item Strategic advise and management control appear as key board/lead tasks.
  \item Country staff includes data collection and support.
\end{itemize}
\textbf{ELI5:} Different roles have different duties, but they all support strategy through control and information.\\
\textbf{Exam write:} ``Connect responsibilities to strategy: boards advise + control; country leaders collect data + execute.''

% ---------------------------
\section*{Slide 14 of 17 (Illustration 1.2: sustained success)}
\textbf{Key idea (from slide):}
\begin{itemize}
  \item Choose 3 companies from 3 different industries (S\&P 500).
\item Give success figures (2023 \& 2013).
  \item Explore possible sources.
\end{itemize}
\textbf{ELI5:} To understand sustained success, compare different companies over time and ask ``why'' they keep winning.\\
\textbf{Exam write:} ``Use an example method: compare industries + success over time + identify sources.''

% ---------------------------
\section*{Slide 15 of 17 (Academic traditions of SM)}
\textbf{Key idea (from slide):}
\begin{itemize}
  \item Economics (Agency theory, Transaction Cost Economics)
  \item Industrial Organization
  \item Strategic Thinking
  \item Organization Studies
  \item Behavioural Sciences (Industrial Psychology)
\end{itemize}
\textbf{ELI5:} Strategy is studied using many fields, not only business.\\
\textbf{Exam write:} ``Explain SM is multi-disciplinary: economics + industry structure + organizations + people behaviour.''

% ---------------------------
\section*{Slide 16 of 17 (SM as academic discipline (1960s, pioneers, ABC Model))}
\textbf{Key idea (from slide):}
\begin{itemize}
  \item SM discipline emerges in the 60s of the 20th century.
  \item Pioneers: Chandler (1962), Andrews (1965), Ansoff (1965).
  \item ABC Model (Academics, Business, Consultants).
  \item Advantages: synergies between academy and practice.
  \item Disadvantages: terminology jungle, lack of coherence.
  \item Boston Consulting Group (1963) appears.
\end{itemize}
\textbf{ELI5:} Strategy science started as a field in the 1960s and mixed academia + practice.\\
\textbf{Exam write:} ``Write SM history: 1960s emergence + pioneers + ABC model synergy/terminology problems.''

% ---------------------------
\section*{Slide 17 of 17 (Bibliography / Further reading)}
\textbf{Key idea (from slide):} Reference list + further reading including ``What is strategy?'' (Porter, 1996).\\
\textbf{ELI5:} These are the sources behind the slides.\\
\textbf{Exam write:} ``Key reading anchor: Porter (1996) about strategy definition and purpose.''

\end{document}

